\section{Decorated integral lattices and short vectors in their duals}
\label{al-section}

We prove three arithmetic inputs used by the coefficient maps of
Section~\ref{am-section}: the existence of the initial decorated
unimodular lattice, its kernel lattice's exact degree, and the existence
of sufficiently short independent vectors in the Hermitian dual.
We also prove the triviality of the actual stabilizers by their
unipotent reductions at two different rational primes.
These arguments use no zero-free region for an $L$-function.
The correspondence with the original construction is recorded
separately in the source audit~\cite{alpo-uniform-source}.

Throughout, $E$ is a CM field, $d=[E:\mathbb Q]=2f$, and
$\mathcal O=\mathcal O_E$ is its maximal order. Write
$\overline{\phantom{x}}$ for complex conjugation,
$\mathfrak D$ for the absolute different,
$\Delta=|\operatorname{disc}E|$, and $\rho=\Delta^{1/d}$.
Let $V$ have dimension $n\geq1$ over $E$ and let $h$ be a
Hermitian form, conjugate-linear in its first argument, positive
definite at every complex embedding. A full lattice means a
finitely generated $\mathcal O$-submodule spanning $V$ over $E$.
For a full lattice $L$, put
\[
 L^\#=\{y\in V:h(x,y)\in\mathcal O\text{ for every }x\in L\}.
\]
It is a full lattice. A lattice is integral if $L\subseteq L^\#$,
and unimodular if equality holds. In particular, integrality concerns
the $E$-valued form, not just the integrality of its rational trace.

\subsection{Residue decorations, indices, and determinant degree}

Let $S$ and $G$ be finite sets of nonzero prime ideals of $\mathcal O$.
The decorations used here consist of a unimodular lattice $N$, a
nonzero $\mathcal O$-linear map
\[
       \ell_P:N\longrightarrow\mathcal O/P\qquad(P\in S),
\]
and an isomorphism
\[
       b_Q:N/QN\longrightarrow(\mathcal O/Q)^n\qquad(Q\in G).
\]
Each map has a unique extension to the localization $N_P$ or $N_Q$.
For example, if $sx\in N$ and $s\notin P$, the extension of
$\ell_P$ is
\begin{equation}\label{al-extension}
                 \ell_P(x)=\bar s^{-1}\ell_P(sx).
\end{equation}
If also $tx\in N$ and $t\notin P$, linearity applied to
$tsx$ proves that this value is independent of $s$.
It proves additivity and linearity on the localization as well.
If a function on all of $V$ is wanted, extend it by zero outside
that localization. No assertion of linearity outside the localization
is made or needed.

\begin{proposition}[An explicit initial decoration]\label{al-base}
For the standard form $h(x,y)=\sum_{j=1}^n\bar x_jy_j$ on $E^n$,
take $N=\mathcal O^n$, let $\ell_P$ be the last coordinate modulo
$P$, and let $b_Q$ reduce all coordinates modulo $Q$.
These choices give decorations for every pair of finite sets $S,G$.
\end{proposition}
\begin{proof}
Testing $h(e_j,y)$ shows that $N^\#\subseteq\mathcal O^n$;
the reverse inclusion follows because conjugation preserves
$\mathcal O$. Thus $N$ is unimodular. The vector $e_n$ proves
that every $\ell_P$ is nonzero. Coordinate reduction is onto and
has kernel $QN$, which proves the assertion for $b_Q$.
Formula~\eqref{al-extension} gives the localized maps, including
their prescribed zero extensions. The construction works at
primes above two and does not require principal prime ideals.
\end{proof}

\begin{lemma}[Duals reverse finite indices]\label{al-dual-index}
If $L_1\subseteq L_2$ are full lattices, then
\[
 (L_1^\#)^\#=L_1,\qquad
 [L_1^\#:L_2^\#]=[L_2:L_1].
\]
Here brackets denote the cardinality of the indicated quotient.
\end{lemma}
\begin{proof}
A full lattice over the Dedekind domain $\mathcal O$ is projective.
The nondegenerate form identifies its Hermitian dual with its
conjugate module dual, so projective biduality proves the first
identity. Explicitly, after localization at a prime, choose a basis
of the free module. The dual basis for the nondegenerate form gives
the original module on taking the dual a second time; equality at
all primes proves the global equality.

For the index identity, localize $L_1\subseteq L_2$ at $P$ and
write it in elementary-divisor form
\[
 (L_2)_P=\bigoplus_j\mathcal O_Pe_j,\qquad
 (L_1)_P=\bigoplus_j\varpi_P^{a_j}\mathcal O_Pe_j,
                 \qquad a_j\geq0.
\]
The corresponding quotient of Hermitian duals at $\bar P$ has
the same length $\sum_ja_j$. This follows by taking the dual
basis: the exponents change from $a_j$ to $-a_j$, and scalar
conjugation changes $P$ to $\bar P$. Since
$N\bar P=NP$, multiplication of the local cardinalities proves
the second identity. The same argument also proves that every
Hermitian dual appearing here is a full lattice.
\end{proof}

\begin{proposition}[The common kernel]\label{al-kernel}
Let $N$ be unimodular and let $\ell_P:N\to\mathcal O/P$
be nonzero for each $P$ in a finite set $S$. Put
\[
 L=\{x\in N:\ell_P(x)=0\text{ for all }P\in S\},
       \qquad t=\prod_{P\in S}NP.
\]
Then $L$ is an integral full lattice and
\begin{equation}\label{al-indices}
 [N:L]=[L^\#:N]=t,\qquad [L^\#:L]=t^2.
\end{equation}
For the archimedean metrics induced by $h$, its normalized ordinary
determinant degree is
\begin{equation}\label{al-degree}
                         \widehat{\deg}L=-\frac{\log t}{d}.
\end{equation}
These conclusions do not require $S$ to be disjoint from $\bar S$.
\end{proposition}
\begin{proof}
A nonzero map to the one-dimensional residue field is onto.
Choose $u_P\in N$ with $\ell_P(u_P)=1$.
The Chinese remainder theorem gives $a_P\in\mathcal O$ such that
$a_P\equiv1\pmod P$ and $a_P\equiv0\pmod Q$ for $Q\ne P$ in $S$.
Thus $\ell_Q(a_Pu_P)$ is $1$ for $Q=P$ and $0$ otherwise.
The map $N\to\prod_{P\in S}\mathcal O/P$ is onto, with kernel
$L$, and hence $[N:L]=t$. Also
$(\prod_{P\in S}P)N\subseteq L\subseteq N$, so $L$ is full.
Since $N$ is integral, so is $L$. Duality gives
$L\subseteq N=N^\#\subseteq L^\#$, and
Lemma~\ref{al-dual-index} gives the remaining indices.

Use ordinary degree as defined in \eqref{am-degree}; the metrics
here are the metrics furnished by $h$ at each embedding.
Hermitian duality negates degree. Indeed it inverts the determinant
ideal up to conjugation and takes the reciprocal determinant norm
at every archimedean place, while conjugation does not change ideal
norm. Since $N=N^\#$, its degree is zero.
Passing to a sublattice of index $t$ changes the determinant ideal
norm by a factor $t$ and leaves the ambient metric unchanged.
This proves \eqref{al-degree}.
\end{proof}

In the actual application $|S|=f$ and $NP=T$ for every $P\in S$.
Consequently
\begin{equation}\label{al-actual-degree}
 [L^\#:L]=T^{2f},\qquad
 \widehat{\deg}L=-\tfrac12\log T,\qquad
 \mu(L)=-\frac{\log T}{2n}.
\end{equation}
The last expression becomes $-\log T/n$ in the doubled slope
convention of the original source. In particular,
\eqref{al-actual-degree} supplies the determinant-degree hypothesis
in Corollary~\ref{am-short} for every decorated representative.
It is not necessary to compute the full structure of $L^\#/L$
to obtain this degree.

\subsection{Unipotent reduction and trivial stabilizers}

\begin{lemma}\label{al-local-torsion}
Let $R$ be a discrete valuation ring of characteristic zero whose
residue field has characteristic $p>0$. If $A\in\mathrm{GL}_n(R)$
has finite order and its reduction is unipotent, then the order
of $A$ is a power of $p$. No condition on the ramification of
$R$ or on whether $p=2$ is needed.
\end{lemma}
\begin{proof}
First suppose $B\equiv I\pmod{\mathfrak m}$ and $B^r=I$ with
$p\nmid r$. Write $X=B-I$. If the entries of $X$ belong to
$\mathfrak m^k$ for $k\geq1$, the binomial identity gives
\[
  0=rX+\sum_{j=2}^r\binom rjX^j.
\]
The sum belongs entrywise to $\mathfrak m^{2k}$ and $r$ is a unit
of $R$, so $X$ belongs entrywise to $\mathfrak m^{2k}$.
Induction and $\bigcap_{k\geq1}\mathfrak m^k=0$ give $B=I$.

Now let the order of $A$ be $p^ar$ with $p\nmid r$ and put
$C=A^{p^a}$. Its order is $r$ and its reduction remains unipotent.
For any unipotent matrix $I+Z$ over a field of characteristic $p$,
$(I+Z)^{p^b}=I+Z^{p^b}=I$ when $p^b\geq n$.
Thus the order of $\bar C$ divides both $r$ and a power of $p$,
so $\bar C=I$. The first paragraph yields $C=I$ and hence $r=1$.
\end{proof}

\begin{lemma}\label{al-finite-isometries}
The isometries of a positive definite Hermitian space over $E$
which preserve a fixed full lattice form a finite group.
The same conclusion holds for the rational stabilizer of any
finite-adelic translate of the completed lattice.
\end{lemma}
\begin{proof}
The real form $B(x,y)=\operatorname{Tr}_{E/\mathbb Q}h(x,y)$
is positive definite on $V\otimes_{\mathbb Q}\mathbb R$.
Choose a $\mathbb Z$-basis of the lattice. Each basis vector can
be sent only to lattice points of its fixed $B$-norm. A discrete
lattice meets a bounded set in finitely many points. The images
of the basis determine the map, proving finiteness.

For the second assertion, let the adelic matrix be $x$ and put
$N_x=V\cap x\widehat N$. Choose a nonzero integer $a$ clearing
the finitely many denominators of both $x$ and $x^{-1}$.
Then $aN\subseteq N_x\subseteq a^{-1}N$, so $N_x$ is a full
lattice. A rational element stabilizing $x\widehat N$ stabilizes
$N_x$, and the first assertion applies. Only the containment
between these two global lattices is needed for finiteness;
no strong approximation statement is being used.
\end{proof}

\begin{proposition}[Two different residue characteristics]
\label{al-two-primes}
Let $\gamma$ be an isometry preserving a full lattice in $V$,
or a rational isometry stabilizing a finite-adelic translate.
Suppose that at places above two different rational primes $p,q$,
the matrices of $\gamma$ are conjugate to integral invertible
matrices with unipotent reductions. Then $\gamma=1$.
\end{proposition}
\begin{proof}
Lemma~\ref{al-finite-isometries} gives finite order.
The embeddings of $E$ in its completions are injective; thus
the order of the matrix is unchanged after these embeddings
and after conjugation. Lemma~\ref{al-local-torsion} makes this
order both a power of $p$ and a power of $q$. It is therefore one.
\end{proof}

This is the actual two-prime argument needed when the level
reduces to an upper triangular unipotent subgroup at the two
selected primes. Merely knowing that a matrix is unipotent after
reduction at one arbitrary prime would not prove the assertion.

There is also a useful single-prime alternative. If $p>dn+1$,
the hypotheses at $p$ alone imply $\gamma=1$.
Indeed a nontrivial finite element of $p$-power order has a power
$\delta$ of order $p$. On the rational vector space of dimension
$dn$, its minimal polynomial divides $X^p-1$, is squarefree, and
is not $X-1$. Thus it is divisible by
$\Phi_p(X)=1+X+\cdots+X^{p-1}$, whose irreducibility follows
from Eisenstein applied to $\Phi_p(X+1)$.
It follows that $p-1\leq dn$, a contradiction.
This recovers the separate large-prime decoration criterion.

If a decorated stabilizer consists of a rational isometry
$\gamma$ together with residue-frame matrices $a_Q$ and
$b_Q(\gamma y)=a_Qb_Q(y)$, then $\gamma=1$ also forces
$a_Q=1$ by surjectivity of $b_Q$. Thus these propositions prove
triviality of the entire decorated stabilizer, not just of its
rational projection.

\subsection{A Gaussian proof of the needed Euclidean transference}

We give the analytic step rather than importing a transference
constant. For a full lattice $\Lambda$ in Euclidean
$\mathbb R^D$, write
\[
 \Lambda^*=\{y:\langle y,x\rangle\in\mathbb Z\ (x\in\Lambda)\},
 \qquad
 \lambda_1(\Lambda^*)=\min_{0\ne y\in\Lambda^*}\|y\|.
\]

\begin{lemma}[Gaussian summation]\label{al-Gaussian-Poisson}
For $s>0$, $u\in\mathbb R^D$, and
$c=\operatorname{covol}(\Lambda)^{-1}$, one has
\begin{equation}\label{al-Gaussian-Poisson-eq}
 \sum_{x\in\Lambda}e^{-\pi s\|x-u\|^2}
 =c s^{-D/2}\sum_{y\in\Lambda^*}
       e^{-\pi\|y\|^2/s}\cos(2\pi\langle y,u\rangle).
\end{equation}
All sums converge absolutely, locally uniformly in $u$.
\end{lemma}
\begin{proof}
A lattice has $O((1+R)^D)$ points in a ball of radius $R$:
disjoint sufficiently small balls about its points lie in a fixed
enlargement of that ball. This proves the asserted convergence
and permits termwise integration on a fundamental parallelepiped.
The Fourier transform, with exponent $-2\pi i\langle y,z\rangle$,
of $e^{-\pi s\|z\|^2}$ is
$s^{-D/2}e^{-\pi\|y\|^2/s}$. In one dimension this follows by
differentiating the transform and integrating by parts, which
gives $G'(y)=-2\pi yG(y)/s$; the Gaussian integral gives
$G(0)=s^{-1/2}$. The latter integral follows by squaring it and
using polar coordinates. Multiplication proves the assertion in
dimension $D$.

Unfolding the periodization over a fundamental parallelepiped
now gives its Fourier coefficients as the terms on the right
of \eqref{al-Gaussian-Poisson-eq}, with the factor $c$.
The Fourier series is absolutely convergent. The difference
between it and the continuous periodization has all Fourier
coefficients zero and hence vanishes: convolution with products
of the one-dimensional Fej\'er kernels gives zero and converges
uniformly to the continuous difference. Pairing $y$ and $-y$
replaces the complex exponentials by the displayed cosines.
\end{proof}

\begin{theorem}[Covering and independent vectors]\label{al-transference}
If every nonzero point of $\Lambda^*$ has norm at least
$\lambda>0$, then every $u\in\mathbb R^D$ has some
$x\in\Lambda$ with
\begin{equation}\label{al-cover}
                         \|u-x\|\leq\frac D{2\lambda}.
\end{equation}
There are $D$ linearly independent vectors of $\Lambda$, each
of norm at most $D/\lambda$.
\end{theorem}
\begin{proof}
For $D=1$, write $\Lambda=a\mathbb Z$ in an orthonormal coordinate.
Then $\lambda\leq1/a$, and rounding to the nearest multiple of
$a$ proves both claims. Dimension zero is immediate.
Suppose $D\geq2$, put $m=D/2\geq1$, and suppose for a
contradiction that no lattice point is within $R=m/\lambda$
of some $u$. Define
\[
 A_s(u)=\sum_{x\in\Lambda}e^{-\pi s\|x-u\|^2},\quad
 Q_s=\sum_{y\in\Lambda^*}e^{-\pi\|y\|^2/s},\quad
 Q_s(u)=\sum_{y\in\Lambda^*}e^{-\pi\|y\|^2/s}
                              \cos(2\pi\langle y,u\rangle).
\]
For $0<\tau\leq1$, the distance assumption and Gaussian summation
give
\begin{align*}
 A_s(u)&\leq e^{-\pi s(1-\tau)R^2}A_{s\tau}(u),\\
 A_{s\tau}(u)&\leq c(s\tau)^{-m}Q_s.
\end{align*}
Indeed the cosines are at most one, and $Q_{s\tau}\leq Q_s$.
Consequently
\begin{equation}\label{al-Gauss-upper}
       Q_s(u)\leq e^{-\pi s(1-\tau)R^2}\tau^{-m}Q_s.
\end{equation}
For $0<\theta\leq1$, the lower bound on the nonzero dual vectors
gives
\[
 Q_s-Q_s(u)
 \leq2e^{-\pi(1-\theta)\lambda^2/s}Q_{s/\theta}.
\]
Also $A_{s/\theta}(0)\leq A_s(0)$ term by term, so Gaussian
summation at zero implies
$Q_{s/\theta}\leq\theta^{-m}Q_s$.
Since $Q_s>0$, combination with \eqref{al-Gauss-upper} yields
\begin{equation}\label{al-Gauss-contradiction}
 1\leq e^{-\pi s(1-\tau)R^2}\tau^{-m}
               +2e^{-\pi(1-\theta)\lambda^2/s}\theta^{-m}.
\end{equation}
Choose
\[
       s=\frac{18\lambda^2}{25m},\qquad
       \tau=\frac49,\qquad\theta=\frac14.
\]
The right side of \eqref{al-Gauss-contradiction} is
\[
 \left(\frac94e^{-2\pi/5}\right)^m
          +2\left(4e^{-25\pi/24}\right)^m<1.
\]
Here is an elementary numerical certificate. The perimeter of
an inscribed regular dodecagon gives
$\pi>3(\sqrt6-\sqrt2)>31/10$.
The first five terms of the exponential series at $31/25$
give $e^{31/25}>17/5$. Also
$e>163/60>19/7$ and
$e^{1/5}>1+1/5+1/50=61/50$, so
\[
 e^{25\pi/24}>e^{16/5}
   >(19/7)^3(61/50)>24.
\]
The two positive bases are therefore smaller than $45/68$ and
$1/6$. Since $m\geq1$, their indicated sum is smaller than
$45/68+1/3=203/204<1$. This contradiction proves
\eqref{al-cover}.

Put $R=D/(2\lambda)$. If the vectors of $\Lambda$ of norm at
most $2R$ failed to span $\mathbb R^D$, let $W$ be their span.
Discreteness gives $\varepsilon>0$ such that every lattice vector
of norm at most $2R+\varepsilon$ already has norm at most $2R$.
Choose $u\in W^\perp$ with $R<\|u\|<R+\varepsilon$.
By \eqref{al-cover} there is $x\in\Lambda$ with
$\|u-x\|\leq R$, and then $\|x\|<2R+\varepsilon$, so
$x\in W$. But orthogonality gives
$\|u-x\|\geq\|u\|>R$, a contradiction.
Select $D$ independent vectors from this spanning finite set.
\end{proof}

\subsection{The codifferent and short Hermitian dual vectors}

The codifferent is
\[
 \mathfrak D^{-1}
  =\{a\in E:\operatorname{Tr}_{E/\mathbb Q}(a\mathcal O)
                                                    \subseteq\mathbb Z\}.
\]
It is invariant under conjugation, because the trace is invariant
under every field automorphism. Its inverse has norm
$N\mathfrak D=\Delta$.

\begin{lemma}[A lower bound in the trace dual]\label{al-trace-dual}
Let $L$ be an integral full lattice and put $Y=L^\#$.
If $0\ne x\in V$ and
\[
       \operatorname{Tr}_{E/\mathbb Q}h(x,y)\in\mathbb Z
                                      \qquad(y\in Y),
\]
then
\begin{equation}\label{al-trace-lower}
             \operatorname{Tr}_{E/\mathbb Q}h(x,x)
                                      \geq\frac d{\rho^2}.
\end{equation}
\end{lemma}
\begin{proof}
Since $Y$ is an $\mathcal O$-module, the assumption implies
$h(x,y)\in\mathfrak D^{-1}$ for every $y\in Y$.
Since $L\subseteq Y$, for $\delta\in\mathfrak D$ and $l\in L$
we obtain
\[
 h(l,\delta x)=\delta\overline{h(x,l)}\in\mathcal O.
\]
Thus $\delta x\in Y$. Applying the previous inclusion with
$y=\delta x$ proves $\delta h(x,x)\in\mathfrak D^{-1}$.
Writing $\beta=h(x,x)$, this says
$\beta\mathfrak D^2\subseteq\mathcal O$.
It is a nonzero integral ideal, so
\[
       |N_{E/\mathbb Q}\beta|\,\Delta^2\geq1.
\]
All $d$ embeddings of $\beta$ are positive real numbers.
Their arithmetic-geometric mean inequality therefore gives
\[
 \operatorname{Tr}_{E/\mathbb Q}\beta
       \geq d\,|N_{E/\mathbb Q}\beta|^{1/d}
       \geq d\Delta^{-2/d}=d\rho^{-2}.
\]
This proves the claim without assuming that $L$ is unimodular
or prescribing the index $[L^\#:L]$.
\end{proof}

\begin{theorem}[Short independent dual vectors]\label{al-short-dual}
Every integral full lattice $L$ admits $n$ linearly independent
vectors $y_1,\ldots,y_n\in L^\#$ over $E$ such that
\begin{equation}\label{al-short-dual-bound}
       \frac1d\operatorname{Tr}_{E/\mathbb Q}h(y_j,y_j)
                                \leq n^2\rho^2\qquad(1\leq j\leq n).
\end{equation}
\end{theorem}
\begin{proof}
Equip $V\otimes_{\mathbb Q}\mathbb R$ with the positive definite
inner product $B(x,y)=\operatorname{Tr}_{E/\mathbb Q}h(x,y)$.
Its real dimension is $dn$, whereas the dimension of $V$ over
$E$ is $n$. The group $Y=L^\#$ is a full Euclidean lattice.
Every vector in its Euclidean dual belongs to $V$: a
$\mathbb Z$-basis of $Y$ has a rational nondegenerate Gram matrix,
whose inverse produces the dual basis inside $V$.
Lemma~\ref{al-trace-dual} therefore gives
\[
                     \lambda_1(Y^*)\geq\sqrt d/\rho.
\]
Theorem~\ref{al-transference} supplies $dn$ rationally independent
vectors of $Y$, each with
\[
 B(y,y)\leq\left(\frac{dn}{\sqrt d/\rho}\right)^2
                                                =dn^2\rho^2.
\]
They span $V$ over $\mathbb Q$, and hence also over $E$.
Select an $E$-basis from them. This proves
\eqref{al-short-dual-bound}. The selected tuple need not be an
$\mathcal O$-basis of $L^\#$; no such assertion is required.
\end{proof}

\begin{corollary}[The arithmetic input for the coefficient maps]
\label{al-coefficient-input}
For each decorated lattice of Proposition~\ref{al-kernel} with
$|S|=f$ and $NP=T$ for every $P\in S$, the degree and short-vector
inputs of Corollary~\ref{am-short} hold with
\[
                  Q=n^2\rho^2=n^2(2v_0)^2,
                         \qquad v_0=\rho/2.
\]
They remain true after any coefficient-field extension with
the induced metrics and normalized degree.
\end{corollary}
\begin{proof}
The degree is \eqref{al-actual-degree}. The mean over all
embeddings of $h(y_j,y_j)$ is its trace divided by $d$,
so Theorem~\ref{al-short-dual} is exactly the mean squared-norm
hypothesis in \eqref{am-short-assumption}. Hermitian pairing
with these dual vectors gives integral independent functionals.
The base-change assertions follow from the norm and archimedean
multiplicity calculation in Section~\ref{am-section}.
\end{proof}

\subsection{Finiteness of the actual decorated classes}

The short-vector theorem also gives a direct finiteness proof for
the class objects used by the coefficient maps. It does not require
a quantitative reduction theory or a class-number estimate.

\begin{theorem}[Finiteness of unimodular Hermitian lattice classes]
\label{al-finite-classes}
Fix the positive Hermitian space $(V,h)$ over $E$.
There are only finitely many isometry classes of full
$\mathcal O$-lattices $N\subset V$ satisfying $N=N^\#$.
\end{theorem}
\begin{proof}
The assertion is immediate when $n=0$, so assume $n\geq1$.
Apply Theorem~\ref{al-short-dual} to $L=N$. Since $N=N^\#$,
it supplies an $E$-basis $y_1,\ldots,y_n$ contained in $N$ with
\[
       \operatorname{Tr}_{E/\mathbb Q}h(y_i,y_i)
                      \leq C,\qquad C=dn^2\rho^2.
\]
Write $G=(h(y_i,y_j))$. Every entry belongs to $\mathcal O$.
At every embedding $\sigma:E\hookrightarrow\mathbb C$,
positivity and Cauchy--Schwarz give
\[
 |\sigma(G_{ij})|^2
       \leq \sigma(G_{ii})\sigma(G_{jj})\leq C^2,
\]
because each positive diagonal value is at most its full trace.
The Minkowski embedding of $\mathcal O$ is a lattice, hence
meets the compact set defined by $|\sigma(a)|\leq C$ in
only finitely many points. For completeness, if
$\omega_1,\ldots,\omega_d$ is an integral basis, the
Minkowski images of the $\omega_i$ form a real basis; its
inverse linear map sends a bounded set to a bounded set,
which contains only finitely many integer coordinate tuples.
Thus only finitely many Gram matrices $G$ can occur.

For each occurring $G$, choose one frame $z_1,\ldots,z_n$
with this Gram matrix, and put
$M=\sum_i\mathcal O z_i$. Any other frame with the same
Gram matrix is carried to this frame by the unique
$E$-linear map sending its $i$th vector to $z_i$; the Gram
identity proves that this map is an isometry. After applying
it, the original lattice therefore satisfies
\[
                         M\subseteq N\subseteq M^\#.
\]
The second inclusion follows from integrality of $N$.
The quotient $M^\#/M$ is finite: in the frame $z_i$, the
invertible matrix $G$ identifies $M^\#$ with a full
fractional lattice in $E^n$, and some nonzero integer
annihilates its quotient by $M$. There are only finitely
many subgroups of a finite abelian group, hence only
finitely many intermediate $\mathcal O$-lattices.
Taking the finite union over the possible $G$ proves the result.
\end{proof}

\begin{corollary}[Finiteness with the residue decorations]
\label{al-finite-decorated}
Fix finite sets $S,G$ of primes of $\mathcal O$. Decorate each
unimodular $N$ by a nonzero $\mathcal O$-linear map
$\ell_P:N\to\mathcal O/P$ for every $P\in S$, and by a
residue frame
\[
           N/QN\simeq(\mathcal O/Q)^n\qquad(Q\in G).
\]
Extend these maps to the appropriate localizations and then
by zero outside them, as above. The decorated objects have
finitely many orbits under the isometry group of $(V,h)$.
They consequently have finitely many orbits under its product
with any subgroup of the finite residue frame groups, acting
in the prescribed way. Their stabilizers in this product are
finite.
\end{corollary}
\begin{proof}
Choose representatives for the finitely many undecorated
classes from Theorem~\ref{al-finite-classes}. For a fixed
representative $N$, choose a finite set of $\mathcal O$-module
generators. A map $N\to\mathcal O/P$ is determined by its
values on these generators, so there are only finitely many
such maps. The module $N/QN$ is finite, so there are only
finitely many residue frames. Their localized extensions are
uniquely determined: if $s\notin P$ and $sx\in N$, the
extended value at $x$ is $\bar s^{-1}\ell_P(sx)$.
Thus adding the extensions introduces no further choices.
There are finitely many decorations of each chosen $N$.
Every decorated object is transported to one of these by an
isometry, which proves finiteness of the isometry orbits.
Passing to a larger group can only merge orbits.

An element stabilizing a decorated object has first component
in the finite group of isometries preserving $N$, by
Lemma~\ref{al-finite-isometries}. The second component belongs
to a finite product of finite matrix groups. The stabilizer
is therefore finite. Its triviality, when required, is the
separate two-prime conclusion of Proposition~\ref{al-two-primes}.
\end{proof}

\subsection{Counterexamples to omitted hypotheses and scope}

The following examples delimit the statements just proved.

\begin{itemize}
\item One residue characteristic is insufficient without a size
condition. For an odd prime $p$, multiplication by $\zeta_p$
on $E=\mathbb Q(\zeta_p)$ preserves $\mathcal O_E$ and the
standard positive Hermitian form, has order $p$, and reduces to
$1$ at the prime $(1-\zeta_p)$. At $p=2$, multiplication by $-1$
on $\mathbb Q(i)$ gives the same obstruction at $(1+i)$.
The two-prime statement retains different rational primes.
\item Positivity cannot be omitted from the finiteness step.
On $\mathbb Q(i)^2$ the indefinite Hermitian form
$h(x,y)=i(\bar x_1y_2-\bar x_2y_1)$ is preserved by the matrices
$\left(\begin{smallmatrix}1&m\\0&1\end{smallmatrix}\right)$
for $m\in\mathbb Z$. They preserve $\mathbb Z[i]^2$,
are unipotent at every prime, and need
not have finite order. The local torsion lemma always assumes
finite order, and the global argument obtains it from positivity.
\item Integrality is essential in Theorem~\ref{al-short-dual}.
Take $E=\mathbb Q(i)$, $h(x,y)=\bar xy$, and
$L=M^{-1}\mathbb Z[i]$ for an integer $M\geq3$.
Then $L^\#=M\mathbb Z[i]$, every nonzero dual vector has mean
squared norm at least $M^2$, whereas $n^2\rho^2=4$.
The omitted hypothesis $L\subseteq L^\#$ fails.
\item The index formula does not imply an arbitrary split
description of the discriminant module. For
$E=\mathbb Q(i)$, $N=\mathbb Z[i]$, and $S=\{P\}$ with
$P=(1+i)$, the kernel is $L=P$ and $L^\#=P^{-1}$.
Thus $L^\#/L\simeq\mathcal O/P^2$, which has cardinality
four but is not $(\mathcal O/P)^2$ as an $\mathcal O$-module:
the latter is killed by $P$ and the former is not.
The source's stronger split description retains the condition
that no selected prime is the conjugate of another.
\item The zero extension of a localized residue map is not
globally additive. In $\mathbb Q(i)$ at $P=(3)$,
both $1/3$ and $2/3$ lie outside $\mathcal O_P$ and have
zero extended value, while their sum $1$ has residue $1$.
Only additivity on the localization is used.
\end{itemize}

These are counterexamples to weakened hypotheses, not to the actual
statements. The arguments above close the standard decorated object's
existence, the kernel lattice's degree, the short-dual input, finiteness
of the decorated classes, and the two-prime stabilizer argument.
They do not identify all automorphic
class lattices with these decorated representatives, construct the
period maps, or prove the remaining global theta estimates.
Those identifications and analytic estimates require their separate
source arguments; no existence package is being substituted for them.


